% =====================================================================
%  TITRE IV : DES VÉHICULES DE PATROUILLE
% =====================================================================
\titrepartie{Titre IV}{Des véhicules de patrouille}

% ---------------------------------------------------------------------
\article{Véhicules de patrouille}

\cl[propriete]{Les véhicules de patrouille sont la propriété exclusive du département. Ils sont mis à la disposition des agents pour la durée d'une garde, en fonction de leur grade et des besoins du terrain.}

\cl{Le dispatcher est seul décisionnaire de l'attribution des véhicules et de la composition des patrouilles.}

\cl{Les agents ne conduisent que les véhicules correspondant à leur grade et aux grades inférieurs.}

\cl[vehicules-division]{L'usage d'un véhicule affecté à une division est interdit aux agents qui n'en sont pas membres confirmés. Seuls le Commanding Officer de la division, son assistant ou le Command Staff peuvent délivrer une dérogation, par écrit.}

\cl[etat-vehicule]{Les véhicules sont restitués propres et en bon état. En cas de dommage survenu au cours de la patrouille, les frais de réparation sont mis à la charge du ou des agents responsables du véhicule, ou pris en charge par les partenaires du département selon les conventions en vigueur.}

\cl{Toute collision impliquant un véhicule de service, quelle qu'en soit la gravité, fait l'objet d'un rapport immédiat et d'une annonce radio.}

\cl[modifications]{Il est interdit d'apporter au véhicule de service une modification esthétique, mécanique ou de plaque d'immatriculation, quelle qu'en soit la nature.}

\sstitre{Dotation embarquée}

\cl[dotation-vehicule]{Chaque véhicule de service est équipé au minimum de :}

\begin{multicols}{2}
\begin{itemize}[leftmargin=5mm,topsep=0pt]
  \item 2 herses
  \item 1 dispositif de dépistage d'alcoolémie
  \item 1 dispositif de dépistage de stupéfiants
  \item 1 plaque de protection balistique
  \item 1 tenue de rechange
  \item 3 fusées éclairantes
  \item 1 extincteur
\end{itemize}
\end{multicols}

\cl{L'agent vérifie la dotation du véhicule à sa prise de service et signale immédiatement tout manquant. À défaut de signalement, le manquant lui est imputé.}

% ---------------------------------------------------------------------
\article{Conduite du véhicule de patrouille}

\cl[permis]{Tout agent doit être titulaire d'un permis de conduire valide dans l'État de San Andreas. Il s'agit d'une condition de recrutement et d'une condition permanente de maintien en service.}

\cl{Hors intervention, les agents sont pleinement soumis au Code de la route. Le non-respect de cette règle expose à des sanctions disciplinaires, indépendamment des suites contraventionnelles.}

\cl[signaux]{En intervention, l'agent active ses signaux lumineux et sonores en fonction du degré d'urgence. Le franchissement d'une intersection s'effectue impérativement signaux lumineux et sonores activés, à vitesse maîtrisée.}

\cl[conduite-raisonnable]{En toute circonstance, l'agent adopte une conduite ne mettant en danger ni sa vie, ni celle des citoyens. Il conserve la maîtrise de son véhicule et adapte sa vitesse à la situation, à la densité de la circulation et à l'état de la chaussée.}

\cl{Les signaux prioritaires n'accordent aucune immunité : un agent qui provoque un accident conserve la responsabilité de sa conduite.}

\cl[transport-tiers]{Il est interdit de transporter des tiers dans un véhicule de patrouille en dehors des cas prévus par une opération, une enquête ou une intervention de police, ou par les dispositions relatives à l'accompagnement civil en patrouille.}

\sstitre{Codes de conduite}

\vspace{2pt}
\noindent
\begin{tabular}{P{22mm} P{126mm}}
\headrow
\thead{Code} & \thead{Signification} \\
\tkey{Code 1} & \tcell{Déplacement sans signal lumineux ni sirène. Respect intégral du Code de la route.} \\
\altrow
\tkey{Code 2} & \tcell{Signaux lumineux activés, sans sirène.} \\
\tkey{Code 3} & \tcell{Signaux lumineux et sirène activés. Déplacement d'urgence.} \\
\altrow
\tkey{Code 4} & \tcell{Fin d'intervention. Situation maîtrisée, aucune assistance nécessaire.} \\
\tkey{Code 5} & \tcell{Anticipation sur une poursuite, surveillance de zone. Les unités en tenue s'abstiennent d'approcher.} \\
\altrow
\tkey{Code 6} & \tcell{Recherche ou vérification sur zone, sans signal lumineux ni sirène.} \\
\tkey{Code 99} & \tcell{\textbf{Agent en danger immédiat.} Toutes les unités disponibles se portent sur les lieux en Code 3. Silence radio général.} \\
\bottomrule
\end{tabular}

\cl{L'emploi d'un code engage l'unité qui l'annonce. L'annonce d'un Code 99 non justifié constitue une faute disciplinaire caractérisée.}

% ---------------------------------------------------------------------
\article{Poursuite véhiculaire}

\chapeau{La poursuite est l'acte de police qui blesse et tue le plus de tiers. Elle n'est jamais une fin en soi : elle est un moyen, et le seul critère d'appréciation est le rapport entre le risque créé et le bénéfice attendu.}

\sstitre{Engagement}

\cl[auto-engagement]{Il est interdit de s'engager de sa propre initiative dans une poursuite en cours, y compris lorsqu'elle passe à proximité immédiate. L'engagement intervient sur demande du dispatcher, de l'unité primaire ou d'un superviseur.}

\cl[nombre-unites]{Le nombre d'unités engagées dans la poursuite est limité en fonction de la gravité du fait poursuivi :}
\begin{itemize}
  \item \stress{délit majeur} : deux unités au maximum ;
  \item \stress{délit mineur} : une seule unité.
\end{itemize}

\clsuite{Ces plafonds ne peuvent être dépassés qu'en cas de nécessité absolue, constatée et annoncée par un superviseur. Les unités non engagées ne suivent pas : elles se positionnent en anticipation, sur les axes, aux points de pose de herses ou en bouclage.}

\sstitre{Répartition des rôles}

\vspace{2pt}
\noindent
\begin{tabular}{P{34mm} P{114mm}}
\headrow
\thead{Rôle} & \thead{Mission} \\
\tkey{Unité primaire} & \tcell{Conduit la poursuite et détient la responsabilité des décisions. Elle est la seule à décider de la poursuite ou de son abandon, sous réserve du pouvoir d'un superviseur.} \\
\altrow
\tkey{Unité secondaire} & \tcell{Assure l'intégralité du compte rendu radio : position, direction, axes empruntés, vitesse, comportement du fuyard. L'unité primaire ne parle pas, elle conduit.} \\
\tkey{Appui aérien} & \tcell{Prend le relais du compte rendu dès qu'il dispose du visuel et se substitue à l'unité secondaire. L'appui aérien est systématiquement sollicité lorsqu'il est disponible.} \\
\altrow
\tkey{Superviseur} & \tcell{Suit la poursuite, autorise les techniques de terminaison, et peut l'interrompre à tout moment. Sa décision d'abandon est sans appel.} \\
\bottomrule
\end{tabular}

\sstitre{Conduite pendant la poursuite}

\cl[contresens]{L'emprunt de la voie de circulation inverse est interdit, sauf nécessité absolue et sur une distance strictement limitée.}

\cl[tirs-fuyard]{Lorsque des coups de feu sont tirés depuis le véhicule poursuivi, les unités décrochent immédiatement et poursuivent à distance de sécurité, hors de portée utile. Aucun tir n'est effectué depuis un véhicule en mouvement.}

\cl[tracking]{Sur décision de l'unité primaire, d'un superviseur ou de l'appui aérien, la poursuite passe en \stress{mode suivi} : l'ensemble des unités repasse en Code 1, décroche à distance et se contente de maintenir le contact visuel ou aérien. Ce mode est employé dès que la poursuite devient plus dangereuse que le fait poursuivi.}

\cl[assistance]{En cas d'accident ou de blessure d'un agent, la \textbf{dernière unité de la file} rompt la poursuite et porte assistance. De même, lorsqu'un objet est jeté du véhicule ou qu'un occupant en descend, la dernière unité se détache pour sécuriser.}

\cl{Les dommages subis par les véhicules sont pris en compte de manière réaliste : un véhicule endommagé est un véhicule ralenti, qui doit être réparé avant reprise du service.}

\sstitre{Techniques de terminaison}

\cl[terminaison]{Aucune technique de terminaison ne peut être mise en œuvre sans demande de l'unité primaire et autorisation d'un superviseur.}

\vspace{4pt}
\noindent
\begin{tabular}{P{30mm} P{118mm}}
\headrow
\thead{Technique} & \thead{Conditions cumulatives} \\
\tkey{Herses} & \tcell{Aucun otage à bord ; chaussée sèche et trafic léger ; emplacement de pose offrant un abri à l'agent ; vitesse du fuyard compatible avec un dérapage maîtrisable ; absence à proximité de station-service, de dénivelé dangereux, d'obstacle mortel ou de groupe de personnes ; gabarit et chargement du véhicule compatibles.} \\
\altrow
\tkey{Contact véhiculaire} & \tcell{Vitesse basse ; accord exprès du superviseur ; interdit sur tout véhicule à deux roues ; interdit à proximité d'un obstacle mortel, d'un dénivelé ou de piétons ; interdit lorsqu'un otage se trouve à bord.} \\
\tkey{Barrage roulant} & \tcell{Accord du superviseur ; bonne visibilité ; trafic léger ; poursuite lente et non violente ; aucun échange de tirs antérieur ; effectifs suffisants. \textbf{Technique de dernier recours.}} \\
\altrow
\tkey{Barrage fixe} & \tcell{Jamais de blocage total à vitesse élevée. Un couloir de dégagement est systématiquement ménagé. Les véhicules composant le barrage sont vides de tout occupant.} \\
\bottomrule
\end{tabular}

\cl[rapport-poursuite]{Toute poursuite, qu'elle aboutisse ou non, donne lieu à un rapport de poursuite dans les vingt-quatre heures. Ce rapport mentionne le fait générateur, les unités engagées, l'itinéraire, la durée, les techniques employées et le motif de la fin de poursuite.}
